\documentclass[12pt,a4paper]{article}

% ==================== 宏包 ====================
\usepackage[UTF8]{ctex}
\usepackage{geometry}
\geometry{left=2.5cm, right=2.5cm, top=2.5cm, bottom=2.5cm}
\setlength{\headheight}{14pt}

\usepackage{booktabs}
\usepackage{longtable}
\usepackage{array}
\usepackage{enumitem}
\usepackage{listings}
\usepackage{xcolor}
\usepackage{hyperref}
\usepackage{fancyhdr}
\usepackage{tabularx}
\usepackage{makecell}
\usepackage{multicol}
\usepackage{tcolorbox}

% ==================== 样式 ====================
\hypersetup{colorlinks=true, linkcolor=black!70, urlcolor=blue!60!black}

\pagestyle{fancy}
\fancyhf{}
\fancyhead[C]{\small 数据结构课程设计\quad 第十一周周报}
\fancyfoot[C]{\thepage}
\renewcommand{\headrulewidth}{0.4pt}
\renewcommand{\footrulewidth}{0.4pt}

\lstset{
    basicstyle=\ttfamily\small,
    keywordstyle=\color{blue}\bfseries,
    commentstyle=\color{green!50!black},
    stringstyle=\color{red!60!black},
    numbers=left,
    numberstyle=\tiny\color{gray},
    frame=single,
    breaklines=true,
    tabsize=4,
    language=SQL
}

\definecolor{headblue}{RGB}{41, 65, 122}
\definecolor{lightgray}{RGB}{245, 245, 245}

\begin{document}

% ==================== 标题 ====================
\begin{center}
    {\LARGE\bfseries\color{headblue} 数据结构课程设计周报}\\[6pt]
    {\Large 第十一周（2026年5月12日\,--\,5月18日）}\\[12pt]
    \begin{tabular}{ll}
        \bfseries 课题名称 & 基于智能体的个性化旅游系统 \\
        \bfseries 指导教师 & 郭岗 \\
        \bfseries 小组成员 & 李佳(2024211216)、王皓轩(2024211213)、崔天睿 \\
    \end{tabular}
\end{center}

\rule{\textwidth}{0.8pt}

% ==================== 一、本周工作概述 ====================
\section*{一、本周工作概述}

本周主要围绕\textbf{系统数据库设计与数据填充}展开工作，完成了数据库表结构设计、真实旅游数据填充、数据库访问层封装、索引优化等任务。数据库采用 SQLite3 作为存储引擎，共设计 12 张业务表，填充了包含 6 个景区/校园地图数据在内的完整测试数据，为用户登录、景点推荐、路线规划、美食查询等功能模块提供了数据基础。

% ==================== 二、数据库设计 ====================
\section*{二、数据库设计与实现}

\subsection*{2.1 总体设计}

数据库采用 SQLite3 嵌入式数据库，以单文件 \texttt{data/tourism.db} 存储全部数据。表结构按照业务实体划分，包含用户系统、景区系统、地图系统、美食系统和社区系统五大模块，共计 12 张表、40 个以上索引：

\begin{center}
\begin{tabular}{l l l}
    \toprule
    \bfseries 模块 & \bfseries 表名 & \bfseries 主要字段 \\
    \midrule
    用户系统 & users & id, username, password, nickname, role, created\_at \\
    & user\_interests & id, user\_id, category, weight \\
    & view\_history & id, user\_id, spot\_id, view\_time \\

    景区系统 & scenic\_spots & id, name, type, category, description, popularity, rating \\
    & buildings & id, area\_id, name, type, total\_floors, has\_elevator \\

    地图系统 & nodes & id, area\_id, name, type, pos\_x, pos\_y, floor \\
    & roads & id, area\_id, from\_node, to\_node, distance, congestion \\
    & indoor\_nodes & id, building\_id, floor, name, type, pos\_x, pos\_y \\
    & indoor\_roads & id, building\_id, from\_node, to\_node, distance \\

    社区系统 & diaries & id, user\_id, title, content, destination\_id, popularity \\
    & ratings & id, user\_id, diary\_id, score \\

    美食系统 & foods & id, area\_id, name, cuisine, restaurant, rating, price \\
    \bottomrule
\end{tabular}
\end{center}

\subsection*{2.2 关键技术选型}

\begin{itemize}[leftmargin=2em]
    \item \textbf{数据库引擎}：SQLite3 3.45.0（amalgamation 方式集成至 C++ 后端）
    \item \textbf{外键约束}：所有跨表关联均定义 \texttt{FOREIGN KEY}，启用 \texttt{PRAGMA foreign\_keys = ON}
    \item \textbf{并发优化}：启用 \texttt{PRAGMA journal\_mode = WAL}，支持读操作与写操作并发执行
    \item \textbf{索引覆盖}：每张表均按查询模式建立组合索引，覆盖分类过滤、排序、关联查询等场景
\end{itemize}

\subsection*{2.3 数据库访问层}

封装了 \texttt{DBConnection} 类（\texttt{db\_connection.h/cpp}），提供统一的数据库操作接口：

\begin{itemize}[leftmargin=2em]
    \item \textbf{自动初始化}：首次连接时检测数据库是否为空，自动执行 \texttt{init\_db.sql} 建表
    \item \textbf{事务支持}：提供 \texttt{begin\_transaction() / commit() / rollback()} 方法
    \item \textbf{查询封装}：模板化查询接口 \texttt{query() / query\_int() / query\_string()}，支持回调式结果处理
    \item \textbf{单例模式}：\texttt{Database::get()} 全局单例，确保整个进程共用同一连接
\end{itemize}

% ==================== 三、数据填充 ====================
\section*{三、真实数据填充}

编写了 \texttt{seed\_data.sql} 数据填充脚本，为各张表补充了完整的测试数据集，具体统计如下：

\begin{center}
\begin{tabular}{l r l}
    \toprule
    \bfseries 表名 & \bfseries 数据量 & \bfseries 说明 \\
    \midrule
    users & 11 条 & 10 名普通用户 + 1 名管理员 \\
    scenic\_spots & 15 条 & 6 个景区 + 9 个校园（含北京著名景点及高校） \\
    nodes & 50 条 & 景区/校园内部导航节点（建筑、入口、路口等） \\
    roads & 62 条 & 节点间连接道路（含距离、拥挤度、理想速度） \\
    foods & 15 条 & 各地美食数据（含菜系、价格、评分） \\
    diaries & 10 条 & 旅游日记（含标题、正文、目的地关联） \\
    ratings & 10 条 & 日记评分数据 \\
    user\_interests & 25 条 & 用户兴趣偏好 \\
    view\_history & 20 条 & 用户浏览历史记录 \\
    \midrule
    \textbf{合计} & \textbf{208+ 条} & \\
    \bottomrule
\end{tabular}
\end{center}

\subsection*{3.1 景区/校园地图数据}

地图数据覆盖 6 个区域，每个区域内包含完整的地图节点和道路网络：

\begin{center}
\begin{tabular}{l r r l}
    \toprule
    \bfseries 区域 & \bfseries 节点数 & \bfseries 道路边数 & \bfseries 特点 \\
    \midrule
    北京邮电大学 & 10 & 15 & 校园主要建筑、南北门互通 \\
    北京大学 & 8 & 10 & 博雅塔、未名湖等景点节点 \\
    清华大学 & 8 & 10 & 水木清华、荷塘月色等景点 \\
    颐和园 & 7 & 8 & 昆明湖、佛香阁、长廊 \\
    故宫博物院 & 9 & 10 & 中轴线宫殿布局 \\
    香山公园 & 8 & 9 & 香炉峰、盘山路径 \\
    \midrule
    \textbf{合计} & \textbf{50} & \textbf{62} & \\
    \bottomrule
\end{tabular}
\end{center}

每条道路包含\textbf{距离（米）}、\textbf{拥挤度（0~1 浮点数）}、\textbf{理想速度（米/秒）}和\textbf{交通工具类型（步行/自行车/电瓶车）}四个权值属性，为 Dijkstra 多策略路径规划提供数据支撑。

\subsection*{3.2 美食数据}

填充了 15 种美食数据，覆盖京菜、川菜、鲁菜、本帮菜、甜点、快餐等 9 个菜系，每项含菜名、餐厅、评分、价格、位置节点关联等信息。其中热门美食（北京烤鸭、宫保鸡丁、水煮鱼）评分 4.7+、热度 8800+。

\subsection*{3.3 社区数据}

10 篇旅游日记覆盖故宫、颐和园、北大、清华、北邮、香山六个目的地，每篇日记包含用户评价和评分数据，为日记推荐与全文检索功能提供数据基础。

% ==================== 四、数据库集成 ====================
\section*{四、后端数据库集成}

\subsection*{4.1 Repository 层设计}

后端采用五层架构，数据库访问集中在 Repository 层，各 Service 模块通过 Repository 接口读取数据库：

\begin{center}
\begin{tabular}{l l l}
    \toprule
    \bfseries Service 模块 & \bfseries 数据库表 & \bfseries 算法依赖 \\
    \midrule
    RecommendService & scenic\_spots, user\_interests, view\_history & MaxHeap + Trie \\
    RouteService & nodes, roads & Dijkstra + TSP \\
    FacilityService & nodes, roads & Dijkstra 距离排序 \\
    DiaryService & diaries, ratings & HashMap + 倒排索引 + Huffman \\
    FoodService & foods, nodes & MaxHeap + EditDistance \\
    \bottomrule
\end{tabular}
\end{center}

\subsection*{4.2 数据库文件位置}

数据库文件统一存放于 \texttt{data/} 目录下：
\begin{itemize}[leftmargin=2em]
    \item \texttt{data/tourism.db} — 运行的 SQLite3 数据库文件（约 200KB）
    \item \texttt{data/init\_db.sql} — 数据库初始化 SQL（12 张表结构 + 索引 + 基础数据）
    \item \texttt{data/seed\_data.sql} — 真实数据填充 SQL（208+ 条记录）
\end{itemize}

% ==================== 五、遇到的问题与解决方案 ====================
\section*{五、遇到的问题与解决方案}

\begin{center}
\begin{tabular}{p{4.5cm} p{7cm}}
    \toprule
    \bfseries 问题 & \bfseries 解决方案 \\
    \midrule
    中文字段存储正确但查询结果显示为乱码 & SQLite3 使用 UTF-8 编码，确保 C++ 后端字符串处理保持一致编码，\texttt{setlocale(LC\_ALL, "")} 设为系统编码 \\
    数据库文件路径硬编码导致部署时找不到文件 & 修改为从配置文件读取 \texttt{db\_path} 参数，main.cpp 传入相对路径 \texttt{data/tourism.db} \\
    景点初始数据只有 5 条，前端信息太少 & 补充至 15 条（6 个景区 + 9 个校园），且 seed\_data.sql 中每张表均有 IGNORE 防重复 \\
    多条道路数据时 Dijkstra 取边未按交通类型过滤 & 在 \texttt{roads} 表中增加 \texttt{transport} 字段，Dijkstra 查询时带 \texttt{WHERE transport=?} 条件 \\
    \bottomrule
\end{tabular}
\end{center}

% ==================== 六、代码量统计 ====================
\section*{六、数据库代码量统计}

\begin{center}
\begin{tabular}{l r}
    \toprule
    \bfseries 文件 & \bfseries 代码行数 \\
    \midrule
    init\_db.sql（表结构 + 索引 + 初始数据） & 301 \\
    seed\_data.sql（真实数据填充） & 244 \\
    db\_connection.h（数据库连接接口） & $\sim$80 \\
    db\_connection.cpp（数据库连接实现） & 245 \\
    \midrule
    \textbf{合计} & \textbf{$\sim$870 行} \\
    \bottomrule
\end{tabular}
\end{center}

% ==================== 七、下周计划 ====================
\section*{七、下周工作计划}

\begin{enumerate}[leftmargin=2em]
    \item \textbf{前端完整联调} — 将 5 个前端页面全部对接真实后端 API，验证数据库数据的端到端流通
    \item \textbf{算法可视化} — MapView 实现 Dijkstra 路径动画绘制，DiaryView 展示 Huffman 压缩率
    \item \textbf{数据增强} — 补充景点数据至 200+、设施数据至 50+、道路边至 200+
    \item \textbf{室内导航} — 实现 buildings / indoor\_nodes / indoor\_roads 表的室内导航功能
    \item \textbf{课程设计报告} — 开始撰写课程设计报告初稿
\end{enumerate}

% ==================== 八、AI辅助开发 ====================
\section*{八、AI辅助开发说明}

本周数据库开发过程中，全程使用 WorkBuddy（AI 编程智能体）辅助完成：
\begin{itemize}[leftmargin=2em]
    \item 数据库表结构设计讨论与 SQL 脚本生成
    \item seed\_data.sql 真实数据填充脚本的编写与调试
    \item DBConnection 数据库访问封装层的代码生成
    \item SQLite3 编码、事务管理等问题排查
    \item 数据统计与周报文档生成
\end{itemize}

\end{document}
